% ============================================================================
% LANGENTWURF LE-010: Gefahren erkennen
% Profil: S (Senioren 65+) | Tier: 2 (Standard)
% ============================================================================

\documentclass[11pt,a4paper]{article}
\usepackage[utf8]{inputenc}
\usepackage[T1]{fontenc}
\usepackage[ngerman]{babel}
\usepackage{geometry}
\usepackage{fancyhdr}
\usepackage{lastpage}
\usepackage{xcolor}
\usepackage{tcolorbox}
\usepackage{enumitem}
\usepackage{longtable}
\usepackage{hyperref}
\usepackage{amssymb}

\geometry{left=2.5cm, right=2.5cm, top=2.5cm, bottom=2.5cm}
\definecolor{fach}{HTML}{b35c00}
\definecolor{gray}{HTML}{666666}
\definecolor{successgreen}{HTML}{2a7a4b}
\definecolor{warnred}{HTML}{c62828}

\tcbuselibrary{skins,breakable}
\newtcolorbox{scorebox}{ colback=white, colframe=fach, fonttitle=\bfseries, title=Didaktik-Score, breakable }
\newtcolorbox{warnbox}[1][]{ colback=warnred!5, colframe=warnred, fonttitle=\bfseries, title=#1, breakable }

\pagestyle{fancy}
\fancyhf{}
\fancyhead[L]{\small\textcolor{gray}{Digitale Grundlagen -- 010 Gefahren erkennen}}
\fancyhead[R]{\small\textcolor{gray}{LE Tier 2 / Profil S}}
\fancyfoot[C]{\small\thepage{} / \pageref{LastPage}}
\renewcommand{\headrulewidth}{0.4pt}

\begin{document}

\begin{titlepage}
    \centering
    \vspace*{2cm}
    {\Large\textcolor{gray}{Digitale Grundlagen -- Senioren 65+}}\\[2cm]
    {\Huge\textcolor{fach}{\textbf{Langentwurf}}}\\[0.5cm]
    {\large Tier 2 | Profil S}\\[2cm]
    {\LARGE\textbf{010 -- Gefahren erkennen}}\\[0.5cm]
    {\large Modul B: Verbindungen \& Sicherheit}\\[2cm]
    \begin{tabular}{rl}
        \textbf{Zeitrahmen:} & 75--90 Minuten \\
        \textbf{Vorkenntnisse:} & E-Mail, Internet-Grundlagen \\
    \end{tabular}
    \vfill
    {\normalsize Stand: \today}
\end{titlepage}

\tableofcontents
\newpage

\section{Bedingungsanalyse}

\subsection{Ausgangssituation}
\begin{itemize}
    \item Senioren sind bevorzugte Zielgruppe für Betrüger
    \item Phishing-Mails werden immer professioneller
    \item Angst vor Betrug, aber Unsicherheit bei der Erkennung
    \item Enkel-Trick, falsche Support-Anrufe, SMS-Betrug
\end{itemize}

\subsection{Typische Situationen}
\begin{itemize}
    \item ``Microsoft hat angerufen, mein Computer sei gehackt''
    \item ``Ich habe eine Mail von der Bank bekommen, sieht echt aus''
    \item ``Mein Enkel braucht dringend Geld -- per WhatsApp''
    \item ``Ich habe auf einen Link geklickt, war das schlimm?''
\end{itemize}

\section{Sachanalyse}

\subsection{Kernkonzepte}
\begin{enumerate}
    \item \textbf{Phishing:} Gefälschte E-Mails/Webseiten, die Daten stehlen wollen
    \item \textbf{Enkel-Trick (digital):} WhatsApp-Nachrichten ``Hallo Oma, neue Nummer''
    \item \textbf{Falsche Support-Anrufe:} ``Microsoft'', ``Apple'', ``Ihre Bank''
    \item \textbf{Warnzeichen:} Druck, Rechtschreibfehler, seltsame Absender
    \item \textbf{Goldene Regel:} Im Zweifel: Nicht klicken, selbst Kontakt aufnehmen
\end{enumerate}

\subsection{Warnzeichen-Checkliste}
\begin{warnbox}[Wann sollten die Alarmglocken läuten?]
\begin{itemize}
    \item Dringlichkeit: ``Sofort handeln'', ``Ihr Konto wird gesperrt''
    \item Geld: ``Überweisen Sie jetzt'', ``Gewinn abholen''
    \item Drohung: ``Anzeige'', ``Inkasso'', ``Polizei''
    \item Geheimhaltung: ``Sagen Sie niemandem davon''
    \item Unbekannter Absender mit bekanntem Namen
    \item Rechtschreib- und Grammatikfehler
    \item Links zu seltsamen Adressen
\end{itemize}
\end{warnbox}

\section{Didaktische Analyse}

\subsection{Sicherheitsrelevanz}
Dies ist die \textbf{wichtigste Sicherheitseinheit}. Finanzieller und emotionaler Schaden durch Betrug ist real.

\subsection{Didaktische Reduktion}
\textbf{Ausgeklammert:} Technische Details (Malware, Viren), komplexe Scam-Varianten.

\textbf{Fokus:} Die häufigsten Betrugsmaschen erkennen und richtig reagieren.

\section{Lernziele}

\begin{tabular}{|c|p{7cm}|c|c|}
\hline
\textbf{Nr.} & \textbf{Die Teilnehmer können...} & \textbf{Stufe} & \textbf{Niveau} \\
\hline
FZ1 & 3 typische Betrugsmaschen benennen & Wissen & Basis \\
FZ2 & Warnzeichen in einer E-Mail erkennen & Anwenden & Basis \\
FZ3 & entscheiden, ob eine Nachricht verdächtig ist & Anwenden & Basis \\
FZ4 & die goldene Regel anwenden (im Zweifel selbst anrufen) & Anwenden & Basis \\
FZ5 & wissen, was zu tun ist, wenn sie auf etwas geklickt haben & Anwenden & Vertiefung \\
\hline
\end{tabular}

\section{Verlaufsplanung}

\begin{longtable}{|p{1cm}|p{2cm}|p{5cm}|p{3cm}|p{2cm}|}
\hline
\textbf{Zeit} & \textbf{Phase} & \textbf{Inhalt} & \textbf{TN-Aktivität} & \textbf{Material} \\
\hline
0--10 & Einstieg & Echte Beispiele zeigen: ``Erkennen Sie den Betrug?'' & Raten, Diskutieren & Beispiel-Mails \\
\hline
10--30 & Phishing & Was ist Phishing? Warnzeichen durchgehen, Beispiele analysieren & Mitmachen & S.1 \\
\hline
30--45 & Enkel-Trick & WhatsApp-Betrug, falsche Nummern, echte Fälle & Zuhören, Fragen & S.2 \\
\hline
45--60 & Support-Anrufe & ``Microsoft ruft nie an!'', Wie reagieren? & Rollenspiel & S.2 \\
\hline
60--75 & Goldene Regeln & Im Zweifel: Selbst anrufen, niemandem erzählen, was tun wenn... & Notieren & S.3 \\
\hline
75--85 & Übungen & Beispiele bewerten, Checkliste, Notfall-Kontakte & Bearbeiten & S.3 \\
\hline
\end{longtable}

\section{Lösungen}

\textbf{Was tun, wenn ich geklickt habe?}
\begin{enumerate}
    \item Ruhe bewahren
    \item Keine Daten eingeben (falls noch nicht geschehen)
    \item Passwort ändern (falls eingegeben)
    \item Bank anrufen (falls Kontodaten betroffen)
    \item Vertrauensperson informieren
    \item Bei Geldverlust: Polizei (110)
\end{enumerate}

\textbf{Goldene Regeln:}
\begin{itemize}
    \item Keine Bank, kein Support ruft an und fragt nach Passwörtern
    \item Im Zweifel: Selbst die offizielle Nummer anrufen
    \item Lieber einmal zu viel misstrauen als einmal zu wenig
    \item Bei Geld-Forderungen: Immer mit Vertrauensperson sprechen
\end{itemize}

\section{Didaktik-Score}

\begin{scorebox}
\begin{tabular}{|c|l|c|c|p{3.5cm}|}
\hline
\# & Dimension & Gew. & Score & Notiz \\
\hline
1 & Differenzierung & 15\% & 8/10 & Alle brauchen Basis \\
2 & Taxonomie & 10\% & 10/10 & Erkennen + Handeln \\
3 & Alltagsrelevanz & 10\% & 10/10 & Höchste Relevanz! \\
4 & Aktivierung & 10\% & 10/10 & Beispiele analysieren \\
5 & Scaffolding & 10\% & 9/10 & Checkliste, Regeln \\
6 & Feedback & 10\% & 9/10 & Beispiele mit Auflösung \\
7 & Sprachliche Klarheit & 10\% & 10/10 & Klar, direkt, ohne Jargon \\
8 & Motivation & 10\% & 10/10 & Schutz motiviert stark \\
9 & Barrierefreiheit & 10\% & 9/10 & Große Beispiele \\
10 & Praktikabilität & 5\% & 8/10 & 85 Min notwendig \\
\hline
\multicolumn{3}{|r|}{\textbf{GESAMT:}} & \textbf{9.4/10} & \textcolor{successgreen}{\textbf{FREIGABE}} \\
\hline
\end{tabular}
\end{scorebox}

\end{document}
